% Generated draft: requires mathematical and source review.
\documentclass[11pt]{article}
\usepackage[margin=25mm]{geometry}
\usepackage{fontspec}
\setmainfont{Latin Modern Roman}
\usepackage{amsmath,amssymb,mathtools,unicode-math}
\setmathfont{Latin Modern Math}
\usepackage{enumitem,needspace,xurl,hyperref}
\setlength{\parindent}{0pt}
\setlength{\parskip}{5pt}
\setlength{\emergencystretch}{3em}
\newcommand{\N}{\mathbb{N}}
\newcommand{\Z}{\mathbb{Z}}
\newcommand{\Q}{\mathbb{Q}}
\newcommand{\R}{\mathbb{R}}
\newcommand{\C}{\mathbb{C}}
\newcommand{\F}{\mathbb{F}}
\newcommand{\A}{\mathbb{A}}
\newcommand{\PP}{\mathbb{P}}
\newcommand{\E}{\mathbb{E}}
\newcommand{\eps}{\varepsilon}
\newcommand{\dd}{\,\mathrm{d}}
\begin{document}
\section*{UnsolvedMath ID: 10900004}
\textbf{Source problem code:} catalogue entry
\subsection{Definitions and mathematical statement}
Let \(E=S^3\setminus\operatorname{int}\nu K\) be the compact exterior of the figure-eight knot \(K\), so \(\partial E\cong T^2\). Taking the compact exterior is necessary to interpret the source's ``complement'' with torus boundary. For an arbitrary homeomorphism \(f:\partial E_1\to\partial E_2\) between two copies, form the closed \(3\)-manifold
\[
 N_f=E_1\cup_f E_2.
\]
Danciger's Problem 1.4 asks whether \(N_f\) admits a convex real-projective structure for every such gluing \(f\).

In the usual properly convex interpretation, this asks for a properly convex domain \(\Omega\subset\R\mathrm P^3\), convex and bounded in some affine chart, and a discrete subgroup \(\Gamma\subset\mathrm{PGL}_4(\R)\) preserving \(\Omega\) and acting freely, properly discontinuously and cocompactly, such that
\[
 N_f\cong\Omega/\Gamma.
\]
The source says convex, not strictly convex; no strict convexity or smooth-boundary hypothesis is imposed. Its precise convex-projective convention should be confirmed against the source.

The supplied source cites Ballas--Danciger--Lee for a positive answer for identity gluing, namely the double of \(E\). The research summary also reports general convex-gluing machinery subject to a holonomy-matching condition, while saying arbitrary gluing maps remain unresolved. The special case and conditional gluing framework do not assert a solution for every \(f\).
\subsection{Short English statement}
Does every closed three-manifold obtained by gluing two figure-eight knot exteriors along their boundary tori admit a convex projective structure? The supplied source reports a positive identity-gluing case and the summary notes conditional gluing tools; the general map remains the stated question.
\subsection{Sources}
\begin{itemize}
\item Local UnsolvedMath registry, permanent ID \(10900004\), exact source code \texttt{catalogue entry}, subsection 1.4, attributed to Danciger; no proposed year supplied.
\item Delp, Hoffoss and Manning, \emph{Problems In Groups, Geometry, and Three-Manifolds} (2015), 1.4, PDF page 1: \url{https://arxiv.org/pdf/1512.04620#page=1}; supplied status-source link \url{https://arxiv.org/abs/1512.04620}.
\item The source cites Ballas--Danciger--Lee for the double of the figure-eight knot exterior. No complete bibliographic entry is supplied. The research summary mentions holonomy-matching convex-gluing machinery without a separate citation.
\item Source-TeX extraction checked against rendered PDF text, accessed 2026-07-29. Source evidence NEEDS\_REVIEW: open in 2015 with no complete modern primary-literature audit recorded. Registry status partially\_solved and classification PARTIAL-PROGRESS retain the identity case and unresolved general gluing question.
\item Dataset attribution: Ulam AI / UnsolvedMath Contributors, CC BY 4.0. Source rights NEEDS\_REVIEW: publicly accessible author-submitted source TeX and PDF, redistribution terms requiring release review.
\end{itemize}
\end{document}
